\documentclass[11pt]{article}
\usepackage[margin=1in]{geometry}
\usepackage{amsmath,amssymb,amsthm,hyperref,enumitem}
\newtheorem{theorem}{Theorem}
\newtheorem{lemma}{Lemma}
\newcommand{\dis}{\operatorname{dis}}
\title{Strictly increasing distortion of odd trigonometric moment curves}
\author{Principia Math In Progress}
\date{21 September 2026}
\begin{document}
\maketitle
\noindent\textbf{Task 1: MathDB 363805.}
For $m\geq2$ put
\[
 \Gamma_m(t)=\frac1{\sqrt m}
 (\cos t,\sin t,\cos3t,\sin3t,\ldots,
                  \cos((2m-1)t),\sin((2m-1)t)).
\]
Both spheres have their geodesic metrics. Thus $\Gamma_m$
is $\gamma_{2k+1}$ from M\'emoli and Smith \cite{MS} with
$m=k+1$. Write $D_m=\dis(\Gamma_m)$. Their Conjecture 6 asks
whether these distortions increase with $m$.

\begin{theorem}\label{main}
For every integer $m\geq2$, $D_{m+1}>D_m$. Moreover,
\[
 D_m\ \nearrow\ 2\pi-r=1.78977584927\ldots,
\]
where $r$ is the unique solution of $\tan r=r$ in $(\pi,3\pi/2)$.
The decimal is illustrative; the limit is specified exactly by $r$.
\end{theorem}

\paragraph{Reduction to an angle.}
The source already gives the kernel and distortion formula
\cite[Proposition 5.5 and Remark 5.8]{MS}:
\begin{equation}\label{kernel}
 K_m(t)=\frac1m\sum_{j=0}^{m-1}\cos((2j+1)t)
       =\frac{\sin(2mt)}{2m\sin t},\qquad
 D_m=\max_{0\leq t\leq\pi}|\arccos K_m(t)-t|,
\end{equation}
with continuous endpoint values $K_m(0)=1$, $K_m(\pi)=-1$.
These identities also follow directly by taking the inner product
of two curve points and summing a geometric series.
For $0<t<\pi/2$, the elementary inequality
$|\sin(mu)|\leq m\sin u$ for $0\leq u\leq\pi$ gives
\[
 K_m(t)=\frac{\sin(2mt)}{2m\sin t}
       \leq\frac{m\sin(2t)}{2m\sin t}=\cos t.
\]
The sine inequality follows by induction from the sine addition
formula and $|\cos u|\leq1$.
Also $K_m(\pi-t)=-K_m(t)$, so the two angle defects are opposite.
Consequently
\begin{equation}\label{defect}
 D_m=\max_{0\leq t\leq\pi/2}F_m(t),\qquad
 F_m(t):=\arccos K_m(t)-t\geq0.
\end{equation}
Both endpoint defects are zero, whereas
$F_m(\pi/(2m))=\pi/2-\pi/(2m)>0$. Every maximizing angle
therefore lies in $(0,\pi/2)$.

\begin{lemma}\label{comparison}
If $m\geq2$ is an integer and $0<t<\pi/2$, then
\begin{equation}\label{strict}
 F_{m+1}\left(\frac{m}{m+1}t\right)-F_m(t)
 \geq c\frac{t}{m+1}>0,
 \qquad c:=1-\frac2{\pi\sqrt3}>\frac12.
\end{equation}
\end{lemma}
\begin{proof}
Fix $s=2mt>0$, set $a=\sin s/s$, and define
\[
 q(u)=\frac{u}{\sin u},\qquad
 g_a(u)=\arccos(aq(u))-u.
\]
We compare $u=t$ with $u'=mt/(m+1)$. The shared phase $s$ gives
$F_m(t)=g_a(t)$ and $F_{m+1}(u')=g_a(u')$.
For $0<u\leq\pi/2$,
\begin{equation}\label{qbounds}
 1\leq q(u)\leq\pi/2,\qquad
 0\leq q'(u)=\frac{\sin u-u\cos u}{\sin^2u}\leq1.
\end{equation}
The bound on $q$ uses $\sin u\geq2u/\pi$.
The numerator of $q'$ starts at zero and has derivative $u\sin u$.
For the upper bound observe that
\[
 \sin^2u-\sin u+u\cos u
 =(u-\sin u)\cos u+
             \sin u(\sin u+\cos u-1)\geq0.
\]
Here $u\geq\sin u$ and $\sin u+\cos u\geq1$ on the stated interval.

If $a\geq0$, then $g_a'(u)\leq-1$ wherever defined. In the
present application $a q(t)=K_m(t)<1$, since the first cosine
in the average is less than one. The monotonicity of $q$ makes
the derivative well-defined throughout $(0,t]$.
If $a<0$, then $s>\pi$, so $|a|\leq1/\pi$ and
$|a|q(u)\leq1/2$. Hence on $(0,\pi/2]$,
\[
 g_a'(u)=-1+\frac{|a|q'(u)}{\sqrt{1-a^2q(u)^2}}
         \leq-1+\frac2{\pi\sqrt3}=-c.
\]
In both cases $g_a'(u)\leq-c$ for $0<u\leq t$.
Integration from $u'$ to $t$ proves \eqref{strict}.
Finally $\pi>3$ implies $3\pi^2>27>16$, which gives $c>1/2$.
\end{proof}

\begin{proof}[Proof of Theorem \ref{main}]
Choose a maximizing angle $t_m\in(0,\pi/2)$ in \eqref{defect}.
Lemma \ref{comparison} gives
\[
 D_{m+1}\geq F_{m+1}\left(\frac{m}{m+1}t_m\right)
            >F_m(t_m)=D_m.
\]
This proves the full monotonicity conjecture.

For the limit, let $b(s)=\sin s/s$ for $s>0$ and $b(0)=1$.
Its global minimum is attained at a unique $r\in(\pi,3\pi/2)$.
Indeed, $b\geq0$ on $[0,\pi]$; on $(\pi,2\pi)$ its derivative
has the sign of $s\cos s-\sin s$, whose derivative is
$-s\sin s>0$. This numerator is negative at $\pi$ and positive
at $3\pi/2$, giving the unique minimum $r$ there.
For $s\geq2\pi$ we have $b(s)\geq-1/(2\pi)$, strictly above
$b(3\pi/2)=-2/(3\pi)$, so no later lobe gives a global minimum.
At $r$ we have $\sin r/r=\cos r$ and therefore
$\arccos b(r)=2\pi-r=:M$.

The proof of the lemma also shows, by letting $u\downarrow0$,
that for any $0<t<\pi/2$ and $s=2mt$,
\[
 F_m(t)=g_{b(s)}(t)<g_{b(s)}(0)=\arccos b(s)\leq M.
\]
Taking a maximizing angle gives $D_m<M$.
Conversely $r/(2m)<\pi/2$ for $m\geq2$, and
\[
 D_m\geq F_m\left(\frac r{2m}\right)
 =\arccos\left(\frac{\sin r}{2m\sin(r/(2m))}\right)
       -\frac r{2m}\longrightarrow\arccos\left(\frac{\sin r}{r}\right)=M.
\]
This proves the limit and completes the theorem.
\end{proof}

\paragraph{Scope, attribution, and verification.}
The target concerns distortion of the embedding $\gamma_{2k+1}$.
It is distinct from distortion of the projection correspondence
$R_{2k+1}$, modulus of discontinuity of the projection, or the
optimal Gromov--Hausdorff distance between the two spheres.
The source's kernel and supremum formula are explicitly attributed;
the contribution is the phase-preserving comparison in
Lemma \ref{comparison} and its consequences.
The limit also explains why a finite plot through about $k=100$
suggested a ceiling near $1.77$; it is not a uniform ceiling.
The minimum of the sinc function is an elementary classical fact,
reproved here for completeness.

The verifier checks the kernel and its odd symmetry as exact
Chebyshev-polynomial identities for $1\leq m\leq30$, and checks
the derivative identities symbolically. Two supplemental 70-digit
groups compare the formulas and the quantitative inequality at
800 inputs, and estimate finite-dimensional maxima and the limit.
Numerical maximization is not an interval certificate and is not
used to prove monotonicity or the limiting value.
All 35 pages of the current target source were read; no target
journal final was identified in the recorded bounded search.
Related-paper read limits are disclosed in the literature audit.
Review is by the authoring agent only.

\begin{enumerate}[label=\textbf{Task \arabic*:},start=2,leftmargin=*]
\item Determine the first two terms of $M-D_m$ with explicit error bounds.
\item Prove uniqueness of the maximizing angle in $(0,\pi/2)$ for every $m\geq2$.
\item Study distortion monotonicity for the normalized even-dimensional
trigonometric moment curve embeddings defined in the source.
\item Determine which nonuniform weights on the odd harmonics preserve
a monotonicity statement as additional harmonics are added.
\end{enumerate}
\small
\begin{thebibliography}{9}
\raggedright
\bibitem{MS} F. M\'emoli and Z. T. Smith,
Embedding-Projection Correspondences for the estimation of the
Gromov-Hausdorff distance, working paper,
\href{https://arxiv.org/abs/2407.03295}{arXiv:2407.03295v1} (2024).
Definition 4 on p.~7; Conjecture 6 on p.~17;
Proposition 5.5 and Remark 5.8 on pp.~18--19.
\end{thebibliography}
\end{document}
